% NB11b -- Mecanica vectorial del Wankel: (a) fuerza de presion sobre un flanco y par en el eje excentrico,
% (b) engranaje de sincronizacion (corona fija r1 = 2e, pinon del rotor r2 = 3e) que impone omega_r = omega_e/3.
% Reemplaza la figura "Torqueexplicado" (dibujada a mano) de Rojas Baron et al. (manuscrito Wankel 4T).
\documentclass[tikz,border=4mm]{standalone}
\usepackage{amsmath}
\usetikzlibrary{arrows.meta,calc,angles,quotes}
\begin{document}
% Estilo comun de los diagramas TikZ del modulo VCR (NB10b) -- paleta del proyecto
\definecolor{manivela}{HTML}{2A78D6}
\definecolor{inferior}{HTML}{4A3AA7}
\definecolor{control}{HTML}{1BAF7A}
\definecolor{superior}{HTML}{EB6834}
\definecolor{tinta}{HTML}{3A3A38}
\definecolor{metal}{HTML}{C3C2B7}
\definecolor{perdida}{HTML}{E34948}
\definecolor{fuerza}{HTML}{C0392B}
\tikzset{
  barra/.style={line width=3.2pt, line cap=round},
  pin/.style={circle, fill=white, draw=tinta, line width=0.9pt, inner sep=1.6pt},
  fijo/.style={circle, fill=tinta, inner sep=2.6pt},
  vec/.style={-{Stealth[length=2.6mm]}, line width=0.9pt},
  F/.style={-{Stealth[length=3mm]}, line width=1.2pt, fuerza},
  cota/.style={{Stealth[length=1.8mm]}-{Stealth[length=1.8mm]}, line width=0.5pt, tinta},
  lbl/.style={font=\small, text=tinta},
}
\newcommand{\tierra}[1]{\fill[pattern=north east lines, pattern color=tinta] ($(#1)+(-5,-3)$) rectangle ($(#1)+(5,-6)$);
  \draw[tinta, line width=0.8pt] ($(#1)+(-5,-3)$) -- ($(#1)+(5,-3)$);}

\definecolor{gasc}{HTML}{E34948}
\begin{tikzpicture}[x=1cm,y=1cm]
  % ---------------- (a) ----------------
  \def\R{3.0}\def\e{0.55}\def\th{40}   % escala de dibujo (K = R/e ~ 5.5 para que se vea la excentricidad)
  \node[lbl, font=\bfseries, anchor=west] at (-4.2,4.3) {(a) Fuerza de presi\'on y par en el eje exc\'entrico};
  % estator (epitrocoide de 2 lobulos)
  \draw[tinta, line width=1.3pt, domain=0:360, samples=200, smooth] plot ({\R*cos(\x)+\e*cos(3*\x)}, {\R*sin(\x)+\e*sin(3*\x)});
  % centro del rotor C (orbita de radio e) ; rotor gira th/3
  \coordinate (O) at (0,0);
  \coordinate (C) at ({\e*cos(\th)},{\e*sin(\th)});
  \pgfmathsetmacro{\ar}{\th/3}
  \coordinate (A1) at ($(C)+({\R*cos(\ar)},{\R*sin(\ar)})$);
  \coordinate (A2) at ($(C)+({\R*cos(\ar+120)},{\R*sin(\ar+120)})$);
  \coordinate (A3) at ($(C)+({\R*cos(\ar+240)},{\R*sin(\ar+240)})$);
  % camara de trabajo entre A1 y A2 (gas)
  \fill[gasc!18] (A1) to[bend left=14] (A2) -- cycle;
  \draw[superior, line width=1.6pt] (A1) to[bend right=10] (A2) to[bend right=10] (A3) to[bend right=10] (A1);
  \node[lbl, gasc, font=\small] at ($(A1)!0.5!(A2)+({0.55*cos(\ar+60)},{0.55*sin(\ar+60)})$) {$p_1$};
  \node[lbl, font=\small] at ($(A2)!0.5!(A3)+({0.6*cos(\ar+180)},{0.6*sin(\ar+180)})$) {$p_2$};
  \node[lbl, font=\small] at ($(A3)!0.5!(A1)+({0.6*cos(\ar+300)},{0.6*sin(\ar+300)})$) {$p_3$};
  % punto medio del flanco 1 y vectores
  \coordinate (M) at ($(A1)!0.5!(A2)$);
  \draw[vec, control] (O) -- (C) node[lbl, pos=0.5, below=4pt, control] {$\vec e$};
  \draw[vec, inferior] (C) -- ($(C)!0.93!(M)$) node[lbl, pos=0.55, left, inferior] {$\vec r_1$};
  % fuerza: normal al flanco, hacia el centro del rotor, linea de accion pasa por C
  \draw[F] (M) -- ($(M)!0.62!(C)$) node[lbl, pos=0.35, right, fuerza] {$\vec F_1=p_1A_f\,\hat n$};
  \draw[fuerza, dashed, line width=0.5pt] ($(M)!1.35!(C)$) -- (M);
  \node[fijo] at (O) {}; \node[pin] at (C) {};
  \node[lbl, below left] at (O) {$O$}; \node[lbl, above right=-1pt] at (C) {$C$};
  \foreach \P/\n in {A1/1,A2/2,A3/3} \fill[tinta] (\P) circle (2.2pt);
  \node[lbl, right] at (A1) {\'apice};
  \draw[manivela!50, dashed] (O) circle (\e);
  \node[lbl, align=left, anchor=north west, font=\small] at (-4.2,-3.4) {%
    El empuje del gas es normal al flanco y su l\'inea\\ de acci\'on pasa por el centro del rotor $C$:\\
    $\vec r_1\parallel\vec F_1\Rightarrow\vec r_1\times\vec F_1=0$, y queda\\[2pt]
    $\tau_{eje}=\displaystyle\sum_{i=1}^{3}\vec e\times\vec F_i,\qquad |\vec F_i|=(p_i-p_{c\acute arter})\,\sqrt3\,R\,b$\\[4pt]
    Verificaci\'on independiente (potencia): $\tau_{eje}\,\omega_e=\sum_i p_i\,\dot V_i$};
  % ---------------- (b) ----------------
  \begin{scope}[shift={(10.2,0.4)}]
    \node[lbl, font=\bfseries, anchor=west] at (-3.0,3.9) {(b) Sincronizaci\'on: $\omega_r=\omega_e/3$};
    \def\ee{0.8}
    \coordinate (O2) at (0,0); \coordinate (C2) at ({\ee*cos(30)},{\ee*sin(30)});
    % corona fija (radio 2e) y pinon del rotor (radio 3e, dentado interior) -- aqui: pinon fijo al bloque, corona en el rotor
    \draw[manivela, line width=1.5pt] (O2) circle ({2*\ee});
    \draw[superior, line width=1.5pt] (C2) circle ({3*\ee});
    \foreach \a in {0,15,...,345} \draw[manivela, line width=0.6pt] ($(O2)+({2*\ee*cos(\a)},{2*\ee*sin(\a)})$) -- ++({0.12*cos(\a)},{0.12*sin(\a)});
    \foreach \a in {0,10,...,350} \draw[superior, line width=0.6pt] ($(C2)+({3*\ee*cos(\a)},{3*\ee*sin(\a)})$) -- ++({-0.12*cos(\a)},{-0.12*sin(\a)});
    \coordinate (P) at ($(O2)+({2*\ee*cos(210)},{2*\ee*sin(210)})$);
    \fill[gasc] (P) circle (2.4pt); \node[lbl, left, gasc, font=\small] at (P) {contacto};
    \node[fijo] at (O2) {}; \node[pin] at (C2) {};
    \draw[cota] (O2) -- (C2) node[lbl, midway, above, font=\footnotesize] {$e$};
    \node[lbl, manivela, font=\small] at (0.3,-2.75) {pi\~n\'on fijo al estator, $r_1=2e$ (azul)};
    \node[lbl, superior, font=\small] at (1.2,3.1) {corona interior del rotor, $r_2=3e$};
    \node[lbl, align=left, anchor=north west, font=\small] at (-3.0,-3.15) {%
      Rodadura sin deslizamiento en el contacto ($v=0$):\\
      $\omega_r\,r_2=\omega_e\,(r_2-r_1)\ \Rightarrow\ \dfrac{\omega_r}{\omega_e}=\dfrac{e}{r_2}=\dfrac13$\\[3pt]
      El eje gira 3 veces por vuelta del rotor; cada c\'amara\\
      completa su ciclo de 4 tiempos en $1080^\circ$ de eje.};
  \end{scope}
\end{tikzpicture}
\end{document}
